\chapter{Language of Category Theory}
\section{Universe}
\begin{definitionenv}
    We call \textbf{universe} any set $\mathcal{U}$ whose elements are sets, satisfying the following conditions:
    \begin{enumerate}[label=(\alph*)]
        \item $\varnothing \in \mathcal{U}$.
        \item If $A\in \mathcal{U}$ and $B\in A$, then $B\in \mathcal{U}$.
        \item If $A\in \mathcal{U}$, then $\mathcal{P}(A)\in \mathcal{U}$.
        \item If $I\in \mathcal{U}$ and $(A_i)_{i\in I}\in \mathcal{U}^{I}$, then $$\bigcup_{i\in I} A_i\in \mathcal{U}.$$
    \end{enumerate}
\end{definitionenv}
\begin{axiomenv}
Every set belongs to some universe.
\end{axiomenv}
\begin{propositionenv}
    Let $\mathcal{U}$ be a universe
    \begin{enumerate}[label=(\arabic*)]
        \item If $A\in \mathcal{U}$ and $B\subseteq A$, then $B\in \mathcal{U}$.
        \item For any $A\in \mathcal{U}$, one has $\{A\}\in \mathcal{U}$.
        \item If $A,B\in \mathcal{U}$, then $\{A,B\}\in \mathcal{U}$.
        \item If $A,B\in \mathcal{U}$, then $(A,B)\in \mathcal{U}$, where $$ (A,B)\coloneq \{\{A\},\{A,B\}\}.$$
        \item If $X,Y\in \mathcal{U}$, then the Cartesian product $X\times Y\in \mathcal{U}$, where $$ X\times Y\coloneq \{(x,y)\mid x\in X, y\in Y\}.$$
        \item If $X,Y\in \mathcal{U}$, then the set of all correspondences from $X$ to $Y$ belongs to $\mathcal{U}$.
        \item If $X,Y\in \mathcal{U}$, then $Y^X\in \mathcal{U}$.
        \item Let $I\in \mathcal{U}$, $(X_i)_{i\in I}\in \mathcal{U}^I$, then $$\prod_{i\in I} X_i \in \mathcal{U}.$$
    \end{enumerate}
\end{propositionenv}
\begin{proofenv}
\quad
\begin{enumerate}[label=(\arabic*)]
    \item By (c), $\mathcal{P}(A)\in \mathcal{U}$ and as $B\in \mathcal{P}(A)$ then from (b), we get $B\in \mathcal{U}$.
    \item By (c), $\mathcal{P}(A)\in \mathcal{U}$, as $\{A\}\subseteq \mathcal{P}(A)$ one can deduce from (1) $\{A\}\in \mathcal{U}$.
    \item By (c) and (a) $$\mathcal{P}(\mathcal{P}(\varnothing))=\{\varnothing,\{\varnothing\}\}\in \mathcal{U}.$$ $A,B\in \mathcal{U}$ then by (2), $\{A\},\{B\}\in \mathcal{U}$, by (d), $\{A,B\}=\{A\}\cup\{B\}\in \mathcal{U}.$
    \item By (2), $\{A\}\in \mathcal{U}$, by (3) $\{A,B\}\in \mathcal{U}$ and again by (3), $\{\{A\},\{A,B\}\}\in \mathcal{U}$.
    \item $X,Y\in\mathcal{U}$, so for $x\in X$, $y\in Y$, by (b), $x,y\in U$. By (4), $(x,y)\in \mathcal{U}$, by (2), $\{(x,y)\}\in \mathcal{U}$ and therefore $$\{x\}\times Y = \bigcup_{y\in Y} \{(x,y)\}\in \mathcal{U}.$$ Still by (2), $\dis X\times Y = \bigcup_{x\in X}\{x\}\times Y\in \mathcal{U}$.
    \item A correspondence from $X$ to $Y$ is of the form $f = ((X,Y), \Gamma)$, where $\Gamma \subseteq X\times Y$. By (5) and (c) and (b) $\Gamma \in \mathcal{U}$. By (4), $f \in \mathcal{U}$, and by (2) $\{(X,Y),\Gamma\}\in \mathcal{U}.$ Finally, by (d) $$\bigcup_{\Gamma \in \mathcal{P}(X\times Y)}\{((X,Y),\Gamma)\} \in \mathcal{U}.$$
    \item $Y^X$ is a subset of the set of correspondences from $X$ to $Y$, so by (1) and (b), $Y^X\in \mathcal{U}$.
    \item Let $\dis X=\bigcup_{i\in I}X_i \in \mathcal{U}$, note that $\dis \prod_{i\in I}X_i \subseteq X^I.$ So by (1) and (7), $\dis \prod_{i\in I}X_i \in \mathcal{U}$.
\end{enumerate}
\end{proofenv}
\begin{remark}
    All we have mentioned are sets. And here the ordered pair $(A,B)$ is given by $(A,B)=\{\{A\},\{A,B\}\}$, we can regard it as a kind of coding method, then we can also regard the tuple $(A,B)$ as a set. This makes the problem easy to discuss.
\end{remark}
\section{Category}
\begin{definitionenv}
    Let $\mathcal{U}$ be a universe. We call $\mathcal{U}$-category $\mathcal{C}$ the data consist of
    \begin{enumerate}[label=(\roman*)]
        \item A set $\Obj(\mathcal{C})$. (The elements of which are called objects.)
        \item For any $(X,Y) \in \Obj(\mathcal{C})\times \Obj(\mathcal{C})$ a set $\mathcal{C}(X,Y)\in \mathcal{U}$. (The elements of which are called morphisms from $X$ to $Y$.)
        \item For any $(X,Y,Z)\in \Obj(\mathcal{C})^3$, we have a composition mapping
$$
            \begin{array}{rcl}
            \mathcal{C}(X,Y)\times \mathcal{C}(Y,Z) &\rightarrow &\mathcal{C}(X,Z)\\ 
            (f,g)&\mapsto & g\circ f
            \end{array}
$$
            satisfying the following axioms:
            \begin{enumerate}[label=(\alph*)]
                \item For any $X\in \Obj(\mathcal{C})$, there exists $\mathrm{Id}_X\in \mathcal{C}(X\times X)$ (called the identity morphism of $X$) such that $Y\in \Obj(\mathcal{C})$, one has
                $$\forall f\in \mathcal{C}(X,Y), f\circ \mathrm{Id}_X = f\ ,\forall g\in \mathcal{C}(Y,X), \mathrm{Id}_X\circ g = g$$
                \item For any $(X,Y,Z,W)\in \Obj(\mathcal{C})^4$, $f\in \mathcal{C}(X,Y)$, $g\in \mathcal{C}(Y,Z)$, $h\in \mathcal{C}(Z,W)$, one has $$ h\circ(g\circ f) = (h\circ g)\circ f.$$
                    If $\Obj(\mathcal{C})\in \mathcal{U}$, then we say that $\mathcal{C}$ is \textbf{$\mathcal{U}$-small}. If $\Obj(\mathcal{C})$ is finite and $\mathcal{C}(X,Y)$ is finite for any $X,Y\in \Obj(\mathcal{C})$, then we say that $\mathcal{C}$ is a \textbf{finite category}.
            \end{enumerate}
    \end{enumerate}
\end{definitionenv}
\begin{exampleenv}
    Let $\mathcal{U}$ be a universe. The elements of $\mathcal{U}$ are sets (by definition). Together with (standard set theoretic) mappings as morphism form a category denote by $\mathcal{U}$-set. The objects of $\mathcal{U}$-set are the elements of $\mathcal{U}$ and for any $A,B\in \mathcal{U}$, 
    $$ \mathcal{U}\text{-}\mathrm{set}(A,B)$$
    is just the set of all mappings from $A$ to $B$.
\end{exampleenv}
\begin{remark}
    In practice, we often emit the mention of the universe $\mathcal{U}$ in which we consider the category theory. With this abbreviation, a $\mathcal{U}$-category is simply called a category. A $\mathcal{U}$-small category is called a small category. Similarly $\mathcal{U}$-set is typically denoted as set.
\end{remark}
\begin{notationenv}[also exmaple]
    \quad
\begin{enumerate}[label=(\arabic*)]
    \item $\mathbf{Top} = \{\text{topological spaces with continuous mappings between them}\}$
    \item $\mathbf{Gp} = \{\text{groups, homomorphisms of monoids}\}$
    
        $\mathbf{Mon} = \{\text{monoid, homomorphisim of monoids}\}$

        $\mathbf{Ab}= \{\text{abelian groups, homomorphisms of groups}\}$
    \item $\mathbf{Ring} = \{\text{Unitary rings, homomorphisms of unitary rings}\}$
    \item $\mathbf{K}\text{-}\mathbf{Mod} = \{\text{left-K-modules, }\mathrm{Hom}_K(M,N)\}$, where $K$ is a unitary ring and $M,N$ left $K$-modules.
    \item $\mathbf{Alg}_K=\{K\text{-algebras, homomorphisms of }K\text{ algebras}\}$, where $K$ is a commutative unitary ring.
        Also, $\mathbf{CAlg}$ for those commutative.
\end{enumerate}
\end{notationenv}
\begin{remark}
    Let $K$ be a commutative unitary ring. Recall that a $K$-algebra is a $K$-module $A$ together with a bilinear mapping $$A\times A\rightarrow A$$ such that $A$ equipped with this bilinear mapping forms a unitary ring.
\end{remark}
\begin{notationenv}
    Let $\mathcal{C}$ be a category, $X,Y\in \Obj(\mathcal{C})$. We often write 
    $$ f: X\rightarrow Y \text{ or } X\overset{f}{\longrightarrow} Y$$
    to denote $f\in \mathcal{C}(X,Y)$.
\end{notationenv}
\begin{propositionenv}
    Let $\mathcal{C}$ be a category, $X\in \Obj(\mathcal{C})$. The identity morphism of $X$ is unique.
\end{propositionenv}
\begin{proofenv}
    $ \Id_X = \Id_X \circ \Id_X' = \Id_X'.$
\end{proofenv}
\begin{definitionenv}
    Let $\mathcal{C}$ be a category and $f:X\rightarrow Y$ be a morphism. We say that $f$ is left invertible if there exists $g:Y\rightarrow X$ such that $g\circ f=\Id_X$. We say that $f$ is right invertible if there exists $h:Y\rightarrow X$ such that $f\circ h=\Id_Y$.
\end{definitionenv}
\begin{propositionenv}
    Suppose that $f$ is left invertible and right invertible. Then $f$ has only one left inverse and only one inverse, and they are equal.
\end{propositionenv}
\begin{proofenv}
    Let $g$ be a left inverse and $h$ be a right inverse of $f$.
    $$g = g\circ \Id_Y = g\circ \left(f\circ h\right)= \left(g\circ f\right)\circ h = \Id_X\circ h = h.$$
\end{proofenv}
\begin{definitionenv}
    Let $\mathcal{C}$ be a category and $f:X\rightarrow Y$ be a morphism in $\mathcal{C}$. If $f$ is left and right invertible, there exists a unique morphism $f^{-1}:Y\rightarrow X$ such that
    $$f\circ f^{-1}=\Id_{Y}, f^{-1}\circ f=\Id_{X}.$$
    In this case, we say that $f$ is invertible (or is an isomorphism) and $f^{-1}$ is called the inverse of $f$.
\end{definitionenv}
\begin{exampleenv}
    $\Obj(\mathcal{C})=\{*\}$, $\mathcal{C}(*,*)=(M,\cdot)$ a monoid.
    $$\begin{array}{rcl}
        \mathcal{C}(*,*)\times\mathcal{C}(*,*)&\rightarrow&\mathcal{C}(*,*)\\
        (a,b)&\mapsto& ab
    \end{array}$$
    identity morphism $\Id_X$ is given by the neutral element of $M$.
\end{exampleenv}
\begin{propositionenv}
    Let $\mathcal{C}$ be a category. $X\overset{f}{\rightarrow}Y\overset{g}{\rightarrow}Z$ be a morphism of $\mathcal{C}$. If $f$ and $g$ are isomorphism, so is $g\circ f$. Moreover,
    $$\left(g\circ f\right)^{-1} = f^{-1}\circ g^{-1}.$$
\end{propositionenv}
\begin{proofenv}
    $$\left(f^{-1}\circ g^{-1}\right) \circ\left(g\circ f\right) = f^{-1} \circ \left(g^{-1}\circ g\right)\circ f = f^{-1} \circ \Id_{Y} \circ f = f^{-1}\circ f = \Id_{X}.$$
    $$\left(g\circ f\right)^{-1} \circ \left(f^{-1}\circ g^{-1} \right) = g\circ\left(f\circ f^{-1}\right)\circ g^{-1} = g \circ \Id_{Y} \circ g^{-1} = g\circ g^{-1} =\Id_Z.$$
\end{proofenv}
\begin{definitionenv}
    Let $\mathcal{C}$ be a category. We denote by $\mathcal{C}^{\op}$ the following category
    \begin{itemize}
        \item $\Obj(\mathcal{C}^{\op}) = \Obj(\mathcal{C})$.
        \item $\forall (X,Y)\in \Obj(\mathcal{C})^2$, $\mathcal{C}^{\op}(X,Y) = \mathcal{C}(Y,X).$
        \item $$\begin{array}{rcl}
            \mathcal{C}^{\op}(X,Y)\times \mathcal{C}^{\op}(Y,Z)&\rightarrow&\mathcal{C}^{op}(X,Z)\\
            (f,g)&\mapsto&f\circ g.
        \end{array}$$
    \end{itemize}
        \begin{center}
    \begin{tikzcd}
X \arrow[r,  "f"] \arrow[dr,  swap,  "f\circ g"] &\displaystyle Y \arrow[d,  "g"] \\
& Z
\end{tikzcd}
\end{center}
\end{definitionenv}
\begin{definitionenv}
    A diagram of morphism in a category $\mathcal{C}$ is a directed graph whose vertices are objects and arrows are morphisms.

    \quad We say that a diagram is commutative if for any pair of vertices $(X,Y)$, the morphism obtained by composing along different paths linking $X$ to $Y$ are the same.
\end{definitionenv}
\begin{exampleenv}
    \quad \newline
    Commutative if $g\circ f =h$.
    \begin{center}
    \begin{tikzcd}
X \arrow[r,  "f"] \arrow[dr,  swap,  "h"] &\displaystyle Y \arrow[d,  "g"] \\
& Z
\end{tikzcd}
\end{center}
Commutative if $g\circ f =k\circ h$.
    \begin{center}
    \begin{tikzcd}
X \arrow[r,  "f"] \arrow[d,  swap,  "h"] &\displaystyle Y \arrow[d,  "g"] \\
Z \arrow[r,"k"]& W
\end{tikzcd}
\end{center}
\end{exampleenv}
\newpage
\section{Functors}
\begin{definitionenv}
    Let $\mathcal{C}$ and $\mathcal{D}$ be categories. We call \textbf{functor} from $\mathcal{C}$ to $\mathcal{D}$ the following data:
    \begin{enumerate}[label=(\alph*)]
        \item A mapping $$\begin{array}{rcl}
            \Obj(\mathcal{C})&\rightarrow&\Obj(\mathcal{D})\\
            X&\mapsto & F(X).
        \end{array}$$
        \item $\forall(X,Y)$ couple of objects in $\mathcal{C}$, a mapping
        $$\begin{array}{rcl}
            \mathcal{C}(X,Y)&\rightarrow&\mathcal{D}(F(X),F(Y))\\
            f&\mapsto & F(f).
        \end{array}$$
        which satisfy the following axioms
        \begin{enumerate}[label=(\roman*)]
            \item $\forall X\in \Obj(\mathcal{C})$, $F(\Id_X)=\Id_{F(X)}$.
            \item For any morphism $X\overset{f}{\rightarrow}Y\overset{g}{\rightarrow}Z$ in $\mathcal{C}$,
            $$F\left(g\circ f\right) = F(g)\circ F(f).$$
        \end{enumerate}
    \end{enumerate}
    We use the expression $F:\mathcal{C}\rightarrow\mathcal{D}$ to denote a functor from $\mathcal{C}$ to $\mathcal{D}$.
\end{definitionenv}
\begin{exampleenv}
    $P:\underline{\mathrm{CRing}}\rightarrow\underline{\mathrm{Set}}$.
    
    $\forall A\in \Obj(\mathbf{CRing})$,
    $$P(A)=\{(x,y,z)\in A^3\mid x^2+y^2=z^2\},$$
    $\forall f:A\rightarrow B$,
    $$ \begin{array}{rrcl}
        P(f):&P(A)&\rightarrow& P(B)\\
        &(x,y,z)&\mapsto&(f(x),f(y),f(z)).
    \end{array}
    $$
\end{exampleenv}
\begin{exampleenv}
    Let $\mathcal{C}$ be a category and $X\in\Obj\left(\mathcal{C}\right)$. We denote by $h_X:\mathcal{C}\rightarrow\underline{\mathrm{Set}}$, sending $Y\in \Obj(\mathcal{C})$ to $\mathcal{C}(X,Y)$.

    If $f:Y\rightarrow Z$ is a morphism in $\mathcal{C}$,
    $$\begin{array}{rrcl}
        h_X(f):& \mathcal{C}(X,Y)&\rightarrow &\mathcal{C}(X,Z)\\
        & g&\mapsto&f\circ g\\
        \Id_{\mathcal{C}(X,Y)}=h_X(\Id_Y):&\mathcal{C}(X,Y)&\rightarrow & \mathcal{C}(X,Y)\\
        & g&\mapsto& \Id_Y\circ g = g
    \end{array}
    $$
        \begin{center}
    \begin{tikzcd}
Y \arrow[r,  "\mu"] \arrow[dr,  swap,  "\nu\circ\mu"] &Z \arrow[d,  "\nu"] \\
& W
\end{tikzcd}
\end{center}
$$\begin{array}{rrcl}
    h_X\left(\nu\circ\mu\right):&\mathcal{C}(X,Y) & \rightarrow &\mathcal{C}(X,W)\\
    & g &\rightarrow &\left(\nu\circ\mu\right)\circ g = \nu\circ \left(\mu\circ g\right)
\end{array}
$$
$$\mu\circ g = h_X(\mu)(g),$$
so,
$$\nu\circ\left(\mu\circ g\right) = h_X(\nu)\left(h_X(\mu)(g)\right),\ h_X(\nu\circ\mu) = h_X(\nu)\circ h_X(\mu).$$
\end{exampleenv}
\begin{exampleenv}
    Let $\mathcal{C}$ be a category. $\Id_\mathcal{C}$ that sends $X\in \Obj(C)$ to $X$ and $f:X\rightarrow Y$ to $f$. This is a functor.
\end{exampleenv}
\begin{exampleenv}
    Let $\mathcal{C}$, $\mathcal{D}$, $\mathcal{E}$ be categories, $F:\mathcal{C}\rightarrow\mathcal{D}$ and $G:\mathcal{D}\rightarrow\mathcal{E}$ be functors. We define a functor
    $$G\circ F:\mathcal{C}\rightarrow \mathcal{E}$$
    such that 
    $$\forall X\in \Obj(\mathcal{C}),\ \left(G\circ F\right)(X)=G(F(X)),$$
    $$\forall f:X\rightarrow Y,\ \left(G\circ F\right)(f) = G(F(f)).$$
\end{exampleenv}
\begin{definitionenv}
    Let $\mathcal{C}$ and $\mathcal{D}$ be categories, and $F:\mathcal{C}\rightarrow \mathcal{D}$ and $G:\mathcal{C}\rightarrow\mathcal{D}$ be functors. We call natural transformation (morphism of function) from $F$ to $G$ any family
    $$\varphi=\left(\varphi_X\right)_{X\in \Obj(\mathcal{C})}$$
    where $\varphi_X:F(X)\rightarrow G(X)$ is a morphism (in $\mathcal{D}$) such that for any morphism $f:X\rightarrow Y$ in $\mathcal{C}$, the following diagram is commutative:
    \begin{center}
    \begin{tikzcd}
F(X) \arrow[r,  "F(f)"] \arrow[d,  swap,  "\varphi_{X}"] &\displaystyle F(Y) \arrow[d,  "\varphi_Y"] \\
G(X) \arrow[r,"G(f)"]& G(Y)
\end{tikzcd}
\end{center}
$$G(f)\circ \varphi_X=\varphi_Y\circ F(f).$$
\end{definitionenv}
\begin{exampleenv}
    Let $\Id_F$ be the family $\left(\Id_{F(X)}\right)_{X\in \Obj(\mathcal{C})}$. This is a morphism of functors from $F$ to $F$, called the identity morphism.
    \begin{center}
    \begin{tikzcd}
F(X) \arrow[r,  "F(f)"] \arrow[d,  swap,  "\Id_{F(X)}"] &\displaystyle F(Y) \arrow[d,  "\Id_{F(Y)}"] \\
F(X) \arrow[r,"F(f)"]& F(Y)
\end{tikzcd}
\end{center}
$$F(f)\circ \Id_{F(X)} = F(f) =\Id_{F(Y)}\circ F(f).$$
\end{exampleenv}
\begin{exampleenv}
    Let $F$, $G$ and $H$ be functors from $\mathcal{C}$ to $\mathcal{D}$, $\varphi: F\rightarrow G$, $\psi:G\rightarrow H$ be natural transformations. We define
    $$\psi\circ \varphi: F\rightarrow H$$
    as the family 
    $$ \left(\psi_X\circ\varphi_X\right)_{X\in \Obj(\mathcal{C})},\ F(X)\overset{\varphi_X}{\longrightarrow}G(X)\overset{\psi_X}{\longrightarrow}H(X).$$
    It is called the composite of $\psi$ and $\varphi$.
\end{exampleenv}
\begin{remark}
    Let $\mathcal{C}$ and $\mathcal{D}$ be categories. Let $\underline{\mathrm{Fun}}(\mathcal{C},\mathcal{D})$ be the data that is composite of
    \begin{itemize}
        \item The set of all functors from $\mathcal{C}$ to $\mathcal{D}$
        \item Natural transformations.
    \end{itemize}
    Then $\underline{\mathrm{Fun}}(\mathcal{C},\mathcal{D})$ forms a category (is some suitable universe). It is called the category of functors from $\mathcal{C}$ to $\mathcal{D}$. If $F$ and $G$ are two functors from $\mathcal{C}$ to $\mathcal{D}$ such that there exists an isomorphism from $F$ to $G$ in $\underline{\mathrm{Fun}}(\mathcal{C},\mathcal{D})$. We say that $F$ and $G$ are isomorphic.
\end{remark}
\begin{propositionenv}
    Let $\mathcal{U}$ be a universe, $\mathcal{C}$ be a $\mathcal{U}$-small category, and $\mathcal{D}$ be a $\mathcal{U}$-category. For any functors $F$ and $G$ from $\mathcal{C}$ to $\mathcal{D}$, the set $\mathrm{Nat}(F,G)$ of natural transformation from $F$ to $G$ belongs to $\mathcal{U}$.
\end{propositionenv}
\begin{proofenv}
    $\mathrm{Nat}(F,G)$ is a subset of $\dis \prod_{X\in \Obj(\mathcal{C})}\mathcal{D}(F(X),G(X))\in \mathcal{U}$.
\end{proofenv}
\begin{notationenv}
    If $\mathcal{C}$ is a category, then $\underline{\mathrm{Fun}}(\mathcal{C},\underline{\mathrm{Set}})$ is denoted as $\hat{\mathcal{C}}$.
\end{notationenv}

\section{Yoneda Lemma}
\begin{theoremenv}
    Let $\mathcal{C}$ be a category, $F:\mathcal{C}\rightarrow \underline{\mathrm{Set}}$ be a functor and $X\in \Obj(\mathcal{C})$. Then
    $$\begin{array}{rrcl}
        \beta:&\mathrm{Nat}(h_X,F)&\rightarrow& F(X)\\
        & \psi&\mapsto&\psi_X(\Id_X)\footnote{$\psi_X:h_X(X)=\mathcal{C}(X,X)\rightarrow F(X)$.}
    \end{array}$$
    is a bijection.
\end{theoremenv}
\begin{proofenv}
    We construct a mapping
    $$\begin{array}{rrcl}
        \alpha:&F(X)&\rightarrow&\mathrm{Nat}(h_X,F)\\
        &x&\mapsto&\alpha(x)=(\alpha(x)_Y)_{Y\in \Obj(\mathcal{C})}
    \end{array}
    $$
    where $\alpha(x)_{Y}:h_X(Y)=\mathcal{C}(X,Y)\rightarrow F(Y)$, $f\mapsto F(f)(x)$. We need to check that $\alpha(x)$ is a natural transformation.
    Namely, for any $ \mu:A\rightarrow B$ in $\mathcal{C}$, we have a commutative diagram
    \begin{center}
    \begin{tikzcd}
\mathcal{C}(X,A) \arrow[d, swap, "\alpha(x)_A"]  \arrow[r, "h_X(\mu)"]& \mathcal{C}(X,B) \arrow[d,  "\alpha(x)_B"] \\
F(A) \arrow[r,"F(\mu)"]& F(B)  
\end{tikzcd}
\end{center}
In fact, 
$\forall f\in \mathcal{C}(X,A)$, $$\alpha(x)_B(h_X(\mu)(f))=\alpha(x)_B(\mu\circ f)=F(\mu\circ f)(x),$$
$$F(\mu)\left(\alpha(x)_A(f)\right) = F(\mu)\left(F(f)(x)\right)=F(\mu\circ f)(x).$$
We then check that $\alpha$ and $\beta$ are inverse of each other.
\begin{itemize}
    \item Let $\psi:h_X\rightarrow F$ be a natural transformation and $x=\psi_X(\Id_X)$. For any $f:X\rightarrow Y$ in $\mathcal{C}$,
    $$\alpha(x)_{Y}(f) = F(f)(x) = F(f)\left(\psi_X(\Id_X)\right)=\psi_Y\left(h_X(f)(\Id_X)\right)=\psi_Y(f).$$
    So, $\alpha(x)=\psi$.
    \item For any $x\in F(x)$,
    $$ \beta\left(\alpha(x)\right)=\alpha(x)_X\left(\Id_X\right)= F(\Id_X)(x) = \Id_{F(X)}(x)=x.$$
\end{itemize}
    \begin{center}
    \begin{tikzcd}
h_X(X) \arrow[r,  "F(f)"] \arrow[d,  swap,  "\psi_X"] &\displaystyle F(Y) \arrow[d,  "\psi_Y"] \\
F(X) \arrow[r,"F(f)"]& F(Y)
\end{tikzcd}
\end{center}
\end{proofenv}
\begin{notationenv}
    Let $\mathcal{C}$ be a category, and $f:Y\rightarrow X$ be a morphism in $\mathcal{C}$. By the theorem, we have a bijection
    $$\begin{array}{rcl}
        \mathrm{Nat}(h_X,h_Y)&\rightarrow &h_Y(X)=\mathcal{C}(Y,X)\\
        \psi&\mapsto&\psi_X(\Id_X).
    \end{array}$$
    We denote by $\upsilon(f):h_X\rightarrow h_Y$ the natural transformation that corresponds to $f$ by this bijection. ($\upsilon(f)= \left(\upsilon(f)_Z\right)_{Z\in \Obj(\mathcal{C})}$)
    $$ \begin{array}{rrcl}
        \upsilon(f)_{Z} :& \mathcal{C}(X,Z)&\rightarrow&\mathcal{C}(Y,Z)\\
        &\mu&\mapsto& h_Y(\mu)(f)=\mu\circ f.
    \end{array}
    $$
\end{notationenv}
\begin{propositionenv}
    $\upsilon: \mathcal{C}^{\op}\rightarrow\hat{\mathcal{C}}$ is a functor.
\end{propositionenv}
\begin{proofenv}
    Let $X\in \Obj(\mathcal{C})$. For any $Z\in \Obj(\mathcal{C})$
    $$\begin{array}{rrcl}
        \Id_{\mathcal{C}(X,Z)}=\upsilon(\Id_{X})_Z:&\mathcal{C}(X,Z)&\rightarrow&\mathcal{C}(X,Z)\\
        &\mu&\mapsto&\mu\circ\Id_X=\mu.
    \end{array}
    $$
    Let $Z\overset{g}{\longrightarrow}Y\overset{f}{\longrightarrow}X$ be morphisms in $\mathcal{C}$.
    $$g\circ_{\mathcal{C}^\op} f=f\circ_{\mathcal{C}} g.$$
    $$\begin{array}{rrcl}
        \upsilon\left(g\circ_{\mathcal{C}^\op} f\right)_{W}= \upsilon\left(f\circ_{\mathcal{C}}g\right)_{W} :& \mathcal{C}(X,W)&\rightarrow &\mathcal{C}(Z,W)\\
        &\mu &\mapsto & \mu\circ \left(f\circ g\right)= \left(\mu\circ f\right)\circ g.
    \end{array}
    $$
    $$\upsilon(g)\left(\upsilon(f)(\mu)\right) = \upsilon(g)\left(\mu\circ f\right) = \left(\mu\circ f\right)\circ g.$$
\end{proofenv}
\begin{propositionenv}
    Let $\mathcal{C}$ be a category and $f:Y\rightarrow X$ be a morphism. Then $f$ is an isomorphism if and only if $\upsilon(f)$ is an isomorphism in $\hat{\mathcal{C}}$.

    Moreover, $$\upsilon(f)^{-1} = \upsilon(f^{-1}).$$
\end{propositionenv}
\begin{proofenv}\quad \newline
    ``$\Rightarrow$'' If $f$ is an isomorphism,
    $$ \upsilon(f)\circ \upsilon(f^{-1}) = \upsilon(f^{-1}\circ f) = \upsilon(\Id_Y) = \Id_{h_Y},$$
    $$ \upsilon(f^{-1})\circ \upsilon(f) = \upsilon(f\circ f^{-1}) = \upsilon(\Id_X) =\Id_{h_X}.$$
    ``$\Leftarrow$'' If $\upsilon(f)$ if invertible and $\varphi = \upsilon(f)^{-1}$. Let $g:y\rightarrow X$ such that $\varphi=\upsilon(g)$.
    $$\upsilon(f\circ g) =\upsilon(g)\circ \upsilon(f) =\Id_{h_X},$$
    $$\upsilon(g\circ f) = \upsilon(f)\circ \upsilon(g) = \Id_{h_Y}.$$
\end{proofenv}